\documentclass[10pt,a4paper]{amsart}
\usepackage{amsmath,amssymb,amsthm,graphicx}
\usepackage{verbatim}
\newenvironment{reference}{\comment}{\endcomment}
%\newenvironment{reference}{}{}
\newenvironment{slogan}{\comment}{\endcomment}
\newenvironment{history}{\comment}{\endcomment}

% For commutative diagrams we use Xy-pic
\usepackage[all]{xy}

% We use 2cell for 2-commutative diagrams.
\xyoption{2cell}
% Theorem environments.
%
\theoremstyle{plain}
\newtheorem{theorem}[subsection]{Theorem}
\newtheorem{proposition}[subsection]{Proposition}
\newtheorem{lemma}[subsection]{Lemma}

\theoremstyle{definition}
\newtheorem{definition}[subsection]{Definition}
\newtheorem{example}[subsection]{Example}
\newtheorem{exercise}[subsection]{Exercise}
\newtheorem{situation}[subsection]{Situation}

\theoremstyle{remark}
\newtheorem{remark}[subsection]{Remark}
\newtheorem{remarks}[subsection]{Remarks}

\numberwithin{equation}{subsection}

% Macros
%
\def\lim{\mathop{\mathrm{lim}}\nolimits}
\def\colim{\mathop{\mathrm{colim}}\nolimits}
\def\Spec{\mathop{\mathrm{Spec}}}
\def\Hom{\mathop{\mathrm{Hom}}\nolimits}
\def\Ext{\mathop{\mathrm{Ext}}\nolimits}
\def\SheafHom{\mathop{\mathcal{H}\!\mathit{om}}\nolimits}
\def\SheafExt{\mathop{\mathcal{E}\!\mathit{xt}}\nolimits}
\def\Sch{\mathit{Sch}}
\def\Mor{\mathop{\mathrm{Mor}}\nolimits}
\def\Ob{\mathop{\mathrm{Ob}}\nolimits}
\def\Sh{\mathop{\mathit{Sh}}\nolimits}
\def\NL{\mathop{N\!L}\nolimits}
\def\CH{\mathop{\mathrm{CH}}\nolimits}
\def\proetale{{pro\text{-}\acute{e}tale}}
\def\etale{{\acute{e}tale}}
\def\QCoh{\mathit{QCoh}}
\def\Ker{\mathop{\mathrm{Ker}}}
\def\Im{\mathop{\mathrm{Im}}}
\def\Coker{\mathop{\mathrm{Coker}}}
\def\Coim{\mathop{\mathrm{Coim}}}

% Boxtimes
%
\DeclareMathSymbol{\boxtimes}{\mathbin}{AMSa}{"02}

%
% Macros for moduli stacks/spaces
%
\def\QCohstack{\mathcal{QC}\!\mathit{oh}}
\def\Cohstack{\mathcal{C}\!\mathit{oh}}
\def\Spacesstack{\mathcal{S}\!\mathit{paces}}
\def\Quotfunctor{\mathrm{Quot}}
\def\Hilbfunctor{\mathrm{Hilb}}
\def\Curvesstack{\mathcal{C}\!\mathit{urves}}
\def\Polarizedstack{\mathcal{P}\!\mathit{olarized}}
\def\Complexesstack{\mathcal{C}\!\mathit{omplexes}}
% \Pic is the operator that assigns to X its picard group, usage \Pic(X)
% \Picardstack_{X/B} denotes the Picard stack of X over B
% \Picardfunctor_{X/B} denotes the Picard functor of X over B
\def\Pic{\mathop{\mathrm{Pic}}\nolimits}
\def\Picardstack{\mathcal{P}\!\mathit{ic}}
\def\Picardfunctor{\mathrm{Pic}}
\def\Deformationcategory{\mathcal{D}\!\mathit{ef}}
\usepackage{cleveref}
\crefname{originaldefinition}{original definition}{original definitions}
\Crefname{originaldefinition}{Original definition}{Original definitions}
\makeatletter
\newlabel{properties-definition-Cohen-Macaulay}{{02IO}{source}{}{sourcepointer-1}{}}
\newlabel{properties-definition-Cohen-Macaulay@cref}{{[originaldefinition][0][]02IO}{[1][0][]source}{}{}{}}
\makeatother
\crefname{originaldefinition}{original definition}{original definitions}
\Crefname{originaldefinition}{Original definition}{Original definitions}
\makeatletter
\newlabel{more-morphisms-definition-lci}{{069F}{source}{}{sourcepointer-2}{}}
\newlabel{more-morphisms-definition-lci@cref}{{[originaldefinition][0][]069F}{[1][0][]source}{}{}{}}
\makeatother
\crefname{originallemma}{original lemma}{original lemmas}
\Crefname{originallemma}{Original lemma}{Original lemmas}
\makeatletter
\newlabel{lemma-fundamental-class-lci}{{0E9Z}{source}{}{sourcepointer-3}{}}
\newlabel{lemma-fundamental-class-lci@cref}{{[originallemma][0][]0E9Z}{[1][0][]source}{}{}{}}
\makeatother
\crefname{originallemma}{original lemma}{original lemmas}
\Crefname{originallemma}{Original lemma}{Original lemmas}
\makeatletter
\newlabel{lemma-shriek-over-CM}{{0E9V}{source}{}{sourcepointer-4}{}}
\newlabel{lemma-shriek-over-CM@cref}{{[originallemma][0][]0E9V}{[1][0][]source}{}{}{}}
\makeatother
\crefname{originallemma}{original lemma}{original lemmas}
\Crefname{originallemma}{Original lemma}{Original lemmas}
\makeatletter
\newlabel{divisors-lemma-depth-2-hartog}{{0E9I}{source}{}{sourcepointer-5}{}}
\newlabel{divisors-lemma-depth-2-hartog@cref}{{[originallemma][0][]0E9I}{[1][0][]source}{}{}{}}
\makeatother
\title{A fundamental class map from the top exterior power of the differentials to the shriek module, an isomorphism on the smooth locus}
\author{The Stacks Project Authors}
\date{stacks-project:duality.tex:local-snapshot-20261012; GNU FDL 1.2 or later}
\begin{document}
\maketitle
\noindent\textit{Editorial note.} All delivered mathematical text is the author's own, copied from the pinned Stacks Project chapter file duality.tex; no statement or proof step was rewritten, reordered or removed.\par\medskip
\noindent\textit{Editorial note.} Preamble provenance: the pinned chapter file opens with an input of the shared preamble and carries no preamble text of its own, so the project's single Stacks-wide preamble staged at sources/stacks-duality/preamble.tex is excerpted verbatim under the recipe fixed by LANE\_BRIEF 22.1, as three spans, lines 8-13 (verbatim package and the reference/slogan/history environments), lines 14-18 (Xy-pic setup) and lines 154-217 (the native theorem declarations and mathematical macros). Nothing was added, removed or redefined; the wrapper class and the amsmath/amssymb/amsthm/graphicx base packages are supplied by the assembler.\par\medskip
\noindent\textit{Editorial note.} The counted statement and proof contain no cross reference, no citation and no included file, so the delivered body is self contained.\par\medskip
\noindent\textit{Editorial note.} Canonical location of the counted proposition: Stacks tag 0EA0 (sources/stacks-duality/pinned-original-source.tex line 7845).\par\medskip
\noindent\textit{Allowed standard prerequisites (editorial).}\begin{itemize}
\item The fundamental class on the local complete intersection locus. Same chapter, Lemma labelled lemma-fundamental-class-lci, cited at proof lines 7894-7900. Provides the map on U.
\item The shriek module is (S\_2) over a Cohen-Macaulay base. Same chapter, Lemma labelled lemma-shriek-over-CM, cited at proof lines 7906-7908. Gives the depth condition needed for the Hartog-type extension.
\item Extending a section across a closed subset of codimension at least two. Divisors, Lemma labelled divisors-lemma-depth-2-hartog, cited at proof lines 7908-7910. Extends the map from U to X.
\end{itemize}\medskip
\ifdefined\hypertarget\hypertarget{sourcepointer-1}{}\fi
\noindent\textit{Source reference.} Original definition 02IO: sources/stacks-chow/tags.txt, source lines 3307--3307. Cross-chapter standard prerequisite located in the project tag list staged in the local snapshot; invoked at proof line None of the counted proof. It is retained only as a navigation pointer; every load-bearing step is inside the delivered span.\par\medskip
\ifdefined\hypertarget\hypertarget{sourcepointer-2}{}\fi
\noindent\textit{Source reference.} Original definition 069F: sources/stacks-chow/tags.txt, source lines 8039--8039. Cross-chapter standard prerequisite located in the project tag list staged in the local snapshot; invoked at proof line None of the counted proof. It is retained only as a navigation pointer; every load-bearing step is inside the delivered span.\par\medskip
\ifdefined\hypertarget\hypertarget{sourcepointer-3}{}\fi
\noindent\textit{Source reference.} Original lemma 0E9Z: sources/stacks-duality/pinned-original-source.tex, source lines 7774--7774. Same-chapter pointer into the author environment lemma at source line 7774; invoked at proof line 7871 of the counted proof. It is retained only as a navigation pointer; every load-bearing step is inside the delivered span.\par\medskip
\ifdefined\hypertarget\hypertarget{sourcepointer-4}{}\fi
\noindent\textit{Source reference.} Original lemma 0E9V: sources/stacks-duality/pinned-original-source.tex, source lines 5935--5935. Same-chapter pointer into the author environment lemma at source line 5935; invoked at proof line 7880 of the counted proof. It is retained only as a navigation pointer; every load-bearing step is inside the delivered span.\par\medskip
\ifdefined\hypertarget\hypertarget{sourcepointer-5}{}\fi
\noindent\textit{Source reference.} Original lemma 0E9I: sources/stacks-chow/tags.txt, source lines 17853--17853. Cross-chapter standard prerequisite located in the project tag list staged in the local snapshot; invoked at proof line 7884 of the counted proof. It is retained only as a navigation pointer; every load-bearing step is inside the delivered span.\par\medskip
\begin{lemma}
\label{lemma-fundamental-class-almost-lci}
Let $f : X \to Y$ be a morphism of schemes. Let $r \geq 0$. Assume
\begin{enumerate}
\item $Y$ is Cohen-Macaulay (Properties, Definition
\ref{properties-definition-Cohen-Macaulay}),
\item $f$ factors as $X \to P \to Y$ where the first morphism is
an immersion and the second is smooth and proper,
\item if $x \in X$ and $\dim(\mathcal{O}_{X, x}) \leq 1$,
then $f$ is Koszul at $x$ (More on Morphisms, Definition
\ref{more-morphisms-definition-lci}), and
\item if $\xi$ is a generic point of an irreducible component of $X$, then
we have
$\text{trdeg}_{\kappa(f(\xi))} \kappa(\xi) = r$.
\end{enumerate}
Then with $\omega_{X/Y} = H^{-r}(f^!\mathcal{O}_Y)$ there is a map
$$
\wedge^r\Omega_{X/Y} \longrightarrow \omega_{X/Y}
$$
which is an isomorphism on the locus where $f$ is smooth.
\end{lemma}

\begin{proof}
Let $U \subset X$ be the open subscheme over which $f$ is a
local complete intersection morphism. Since $f$ has relative
dimension $r$ at all generic points by assumption (4) we
see that the locally constant function of
Lemma \ref{lemma-fundamental-class-lci}
is constant with value $r$ and we obtain a map
$$
\wedge^r\Omega_{X/Y}|_U = \wedge^r \Omega_{U/Y}
\longrightarrow
\omega_{U/Y} = \omega_{X/Y}|_U
$$
which is an isomorphism in the smooth points of $f$ (this locus
is contained in $U$ because a smooth morphism is a local complete
intersection morphism). By Lemma \ref{lemma-shriek-over-CM}
and the assumption that $Y$ is Cohen-Macaulay
the module $\omega_{X/Y}$ is $(S_2)$.
Since $U$ contains all the points of codimension $1$ by condition (3)
and using Divisors, Lemma \ref{divisors-lemma-depth-2-hartog}
we see that $j_*\omega_{U/Y} = \omega_{X/Y}$.
Hence the map over $U$ extends to $X$ and the proof
is complete.
\end{proof}

\end{document}
